% =====================================================================
% Final Defense Compliance Report
% Page numbers are read from the manuscript's .aux file through the xr
% package, so compile Balbin_Archie_W3.tex first; the numbers here then
% always match the current manuscript.
% =====================================================================
\documentclass[12pt,letterpaper]{article}
\usepackage[top=1in,bottom=1in,left=1.25in,right=1in]{geometry}
\usepackage[T1]{fontenc}
\usepackage{helvet}
\renewcommand{\familydefault}{\sfdefault}
\usepackage{array}
\usepackage{longtable}
\usepackage{fancyhdr}
\usepackage{xr}
\externaldocument{Balbin_Archie_W3}

\newcommand{\sysname}{QueuEx}

\pagestyle{fancy}
\fancyhf{}
\fancyhead[R]{\thepage}
\renewcommand{\headrulewidth}{0pt}
\setlength{\parindent}{0pt}
\renewcommand{\arraystretch}{1.25}

\newcolumntype{S}{>{\raggedright\arraybackslash}m{5.55cm}}
\newcolumntype{A}{>{\raggedright\arraybackslash}m{5.55cm}}
\newcolumntype{P}{>{\centering\arraybackslash}m{2.3cm}}
\newcommand{\pg}[1]{\large #1}

% One table head, reused by every section.
\newcommand{\complianceHead}{%
\hline
\multicolumn{1}{|c|}{\bfseries SUGGESTIONS/} &
\multicolumn{1}{c|}{\bfseries ACTION} &
\multicolumn{1}{c|}{\bfseries PAGE} \\
\multicolumn{1}{|c|}{\bfseries RECOMMENDATIONS} &
\multicolumn{1}{c|}{\bfseries TAKEN} &
\multicolumn{1}{c|}{\bfseries NUMBER} \\
\hline}

\begin{document}

\begin{center}
  {\large\bfseries FINAL DEFENSE COMPLIANCE REPORT}
\end{center}

\vspace{0.4em}
\noindent{\bfseries TITLE: QUEUEX: A COMPUTER VISION-BASED SMART QUEUE
MANAGEMENT AND PREDICTIVE SERVICE DELIVERY SYSTEM}

\vspace{1em}
\begin{tabular}{@{}p{3.6cm}p{0.6cm}l@{}}
RESEARCHERS   & : & ARCHIE D. BALBIN \\
              &   & JOSHUA ROVIC P. SANOTA \\[0.8em]
ADVISER       & : & FRITZGERALD ENRIC T. IMPERIAL, MIT \\[0.8em]
CHAIRMAN      & : & ANNA LORETTA C. ROMULO, PhD \\
PANEL MEMBERS & : & RONEL B. SIMON, DIT \\
              &   & EMMANUEL N. ENALPE III \\[0.8em]
SECRETARY     & : & LIELA B. GARDOSE, MBA \\
\end{tabular}

\vspace{1.2em}
\noindent This report lists the panel's suggestions from the pre-final
defense on 29~September~2026 and the action taken on each. Part~II
covers the review of the system by the school's Technology Resources
Center on 10~November~2026, and Part~III lists the other changes made to
the system since the pre-final defense. Page numbers refer to the
revised manuscript.

\vspace{1em}
\noindent{\bfseries PART I. PRE-FINAL DEFENSE SUGGESTIONS}

\vspace{0.5em}
\small
\begin{longtable}{|S|A|P|}
\complianceHead
\endfirsthead
\complianceHead
\endhead
\hline
\endfoot

\textbf{Preliminaries:}\newline
Wrong content for the Editor's Certification.
&
Replaced the copied certification with a statement of editorial review
of grammar, spelling, punctuation, sentence structure, clarity, and
consistency, with the editor's name above the signature line.
&
\pg{\pageref{cr:editor}} \\ \hline

\textbf{Background of the Study:}\newline
Introduce queue management in your own words.
&
Added an original opening paragraph explaining what a queue management
system does, how it orders service and shows queue status, and how it
helps students and staff at a college cashiering office during busy
periods.
&
\pg{\pageref{cr:background}} \\ \hline

\textbf{Research Instrument (p.~23):}\newline
The operational log is for System Performance.
&
Moved the description of the operational log and observation record
from the Research Instrument section to System Performance, which it
serves. It identifies the operational log as the record the system
generates of detection, recognition, queue-number issuance, ticketing,
prediction, and dashboard actions, the source of the live recognition
and issuance results, kept separate from the manual observation record
used for the wait-time results. The Research Instrument section now
points to it.
&
\pg{\pageref{cr:oplog}} \\ \hline

\textbf{Reliability (p.~40):}\newline
How was Cronbach's alpha computed? Treat the respondents as one group,
not separately as staff and students.
&
The user-evaluation alpha is now computed on all respondents as one
group, students and staff pooled, on the six statements both groups
answered (two each for usability, reliability, and clarity). The item
variances, the variance of the total scores, and the resulting
coefficient are shown step by step with the formula. The user
acceptance test's alpha is likewise computed on all its respondents as
one group, with its worked calculation. The baseline questionnaire is
likewise computed on all ninety respondents (83 students and 7 staff) as
one group, by section and as a whole, with its worked calculation.
&
\pg{\pageref{tab:sop1-reliability}, \pageref{tab:sop4-reliability},
\pageref{tab:uat-reliability}} \\ \hline

\textbf{Appendix (p.~84):}\newline
Present the questionnaire using the actual form.
&
The baseline appendix now shows the actual student Google Form as
nineteen screenshots, in the order respondents saw it. The staff form
is reproduced from its content.
&
\pg{\pageref{app:sop1-questionnaire}} \\ \hline

\textbf{System Evaluation:}\newline
Provide a statistician's report on the system-evaluation metrics.
&
A user acceptance test was conducted with 79 students in two sessions,
on 16 and 17~November~2026, conducted by the BS Computer Science
researchers together with a BS Information Systems student, using the
questionnaire added as an appendix. Students rated the
completed system Strongly Agree (4.25, Cronbach's alpha 0.941). In the
second session, with the camera running, all 19 face registrations were
saved and the camera issued 10 numbers, a median of 2.2~s after it
could. The results are reported in a new results section. \textbf{Still
pending:} a UAT with cashier staff at the dashboard, and the
statistician's report, which will be attached to the manuscript.
&
\pg{\pageref{cr:uat}, \pageref{sec:uat-results}, \pageref{app:uat}} \\ \hline

\textbf{System:}\newline
Show that the system can scale to the whole university (10,000
students).
&
Replaced the face matcher, which reloaded every enrolled face from the
database for each person, with an in-memory index that scores a face
against all enrolled students at once. At 10,000 enrolled students a
match now takes about 3~ms instead of 10~s and uses about 20~MB. The
backend was load-tested again with up to 400 simultaneous users; no
request received an HTTP error.
&
\pg{\pageref{subsec:load-test}--\pageref{tab:scale-load}} \\ \hline

\textbf{System:}\newline
Show a data-privacy consent screen before registration.
&
The mobile application shows a Data Privacy Notice and Consent screen
before any account is created. It states what is collected, why, who
may access it, how long it is kept, and the student's rights under
R.A.~10173. No account is created until the student taps
\emph{I agree}; the backend records the version agreed to and refuses
to create an account, enrol a face, or accept a request for a number
without it. Existing accounts must accept at their next sign-in.
&
\pg{\pageref{cr:consent}} \\ \hline

\textbf{System:}\newline
Define a data-retention policy for student data.
&
An account with no sign-in or queue request for one year is archived:
deactivated, with its face data erased. An account archived for one
more year is deleted. A daily background process applies the policy,
staff can archive or delete accounts at any time from the dashboard's
Data Retention page (for example, graduates), and students can delete
their own account from the app.
&
\pg{\pageref{cr:retention}} \\ \hline

\end{longtable}

\normalsize
\vspace{1em}
\noindent{\bfseries PART II. REVIEW BY THE TECHNOLOGY RESOURCES CENTER
(10 NOVEMBER 2026)}

\vspace{0.5em}
\noindent The school's Technology Resources Center (TRC) reviewed a
demonstration of the system. With the adviser, the researchers adopted
the suggestions that keep the panel-approved design, and recorded the
others in the scope and delimitation and the recommendations.

\vspace{0.5em}
\small
\begin{longtable}{|S|A|P|}
\complianceHead
\endfirsthead
\complianceHead
\endhead
\hline
\endfoot

\textbf{Payment form:}\newline
The student should state the kind of payment and the amount, and the
cashier should see it.
&
\textbf{Adopted.} The request for a number now includes a payment-slip
form laid out like the cashier's paper Payment Slip (Tuition, Books,
Special Exam with the subject, Others with a description), with an
amount per line and a TOTAL. The cashier sees the lines and the TOTAL on
the dashboard when the student is called.
&
\pg{\pageref{cr:presentation}} \\ \hline

\textbf{Mobile application:}\newline
Improve the terminology, user interface, and flow so users are not
confused.
&
\textbf{Adopted.} One term is used for each step (\emph{Get a Queue
Number}, \emph{Request Number}, \emph{Face check}, \emph{Your number}),
and \emph{window} is used throughout. The app shows a three-step
progress indicator, who is being served at each window, a warning before
a long wait, an \emph{It's your turn} screen with the window number, and
a notice when a number is skipped. A \sysname{} loading screen was also
added at the researchers' initiative.
&
\pg{\pageref{cr:presentation}} \\ \hline

\textbf{Computer vision:}\newline
The face check at the camera does not shorten the line, since students
still queue at the camera; issue the number from the app instead.
&
\textbf{Not adopted; stated as a delimitation.} Identifying the student
by camera before a number is issued is the computer-vision component
required of this Computer Science capstone and approved by the panel,
and the adviser decided to keep it. The point is stated in the scope and
delimitation, and an app-issued number with a passive camera is listed
as a recommendation for future work.
&
\pg{\pageref{cr:trc-scope}, \pageref{sec:recommendations}} \\ \hline

\textbf{Other offices:}\newline
A document-request organizer for the Registrar.
&
\textbf{Outside the scope}, stated in the scope and delimitation.
&
\pg{\pageref{cr:trc-scope}} \\ \hline

\end{longtable}

\normalsize
\vspace{1em}
\noindent{\bfseries PART III. OTHER CHANGES TO THE SYSTEM SINCE THE
PRE-FINAL DEFENSE}

\vspace{0.5em}
\small
\begin{longtable}{|>{\raggedright\arraybackslash}m{4.2cm}|>{\raggedright\arraybackslash}m{6.9cm}|P|}
\hline
\multicolumn{1}{|c|}{\bfseries CHANGE} &
\multicolumn{1}{c|}{\bfseries DESCRIPTION} &
\multicolumn{1}{c|}{\bfseries PAGE} \\ \hline
\endfirsthead
\hline
\multicolumn{1}{|c|}{\bfseries CHANGE} &
\multicolumn{1}{c|}{\bfseries DESCRIPTION} &
\multicolumn{1}{c|}{\bfseries PAGE} \\ \hline
\endhead
\hline
\endfoot

Calling by window
&
Each staff account takes one window and presses \emph{Call next} to call
the next number to it; all staff see the same waiting list. A switch on
the dashboard changes to automatic calling.
&
\pg{\pageref{cr:calling}} \\ \hline

Faster confirmation
&
A person whose face has been captured is confirmed after one second in
the queue zone instead of after a fixed number of camera updates, which
had become slow over the internet to the cloud backend.
&
\pg{\pageref{subsec:candidate}} \\ \hline

Wait estimate
&
The student's wait is now estimated by simulating the line window by
window; every issued number records its estimate and its real wait, and
the dashboard scores the estimate as numbers are served.
&
\pg{\pageref{subsubsec:queue-engine}} \\ \hline

Algorithms
&
Grouped into four major algorithms --- person detection and tracking,
presence confirmation, face-based identity validation and number
issuance, and wait-time estimation --- with calling and ticket
generation stated as supporting procedures. The four are stated with
their steps in the Methodology (Section~2.8.4), and their equations and
parameter values are in an appendix.
&
\pg{\pageref{cr:algorithms}, \pageref{app:algorithms}} \\ \hline

Wait-time model
&
The Gradient Boosting predictor is stated as trained and evaluated but
not deployed, with the reasons, in the scope and delimitation; the
deployed estimate simulates the line, with the people-ahead regression
as fallback. Retraining on the system's own wait log is recommended.
&
\pg{\pageref{cr:gb-scope}, \pageref{sec:ml-training-protocol}} \\ \hline

Diagrams
&
The conceptual paradigm, the system architecture, and the
detection-to-ticket pipeline were redrawn for the current system:
payment slip, consent, calling by window, the line-simulation wait
estimate, the retention service, and the current API routers.
&
\pg{\pageref{fig:conceptual-framework}, \pageref{fig:system-architecture}, \pageref{fig:detection-ticket-pipeline}} \\ \hline

Access control
&
Every dashboard page except the public display now requires a staff
sign-in, and the camera stream and the full queue list are staff-only.
&
\pg{\pageref{cr:presentation}} \\ \hline

Fixes between the UAT sessions
&
Face photos are processed off the server's request loop, so many
students registering at once no longer stall everyone else; face
sign-in opens only the account last used on that phone, after a student
saw a classmate's name; the app's errors say what to do; the camera
tracks at full speed with several people in view; and a five-minute
heads-up before the turn, suggested by a student, was added. The second
session lost no registration and had no wrong sign-in.
&
\pg{\pageref{subsec:uat-fixes}} \\ \hline

\end{longtable}

\normalsize
\vspace{2em}
\noindent Prepared by:

\vspace{2.5em}
\begin{tabular}{@{}p{7cm}p{7cm}@{}}
\rule{6cm}{0.4pt} & \rule{6cm}{0.4pt} \\
ARCHIE D. BALBIN & JOSHUA ROVIC P. SANOTA \\
Researcher & Researcher \\
\end{tabular}

\vspace{2em}
\noindent Noted by:

\vspace{2.5em}
\noindent\rule{6cm}{0.4pt}\\
FRITZGERALD ENRIC T. IMPERIAL, MIT\\
Adviser

\end{document}
